% Pyramide, Kegel, Kugel – bank/pyramide-kegel-kugel/ (zusammenbau v0.4, Heft)
\einheitenkopf[t3]{Pyramide, Kegel, Kugel}
\setcounter{aufgabe}{20}

% Anlauf (ohne Prüfkennung)
\begin{aufgabe}{Pyramide – Anlauf}
\begin{teile}
\teil Strecke von der Spitze senkrecht zur Mitte der Grundfläche – Höhe, Seitenhöhe oder Seitenkante? Kreuze an.\\ \kreuz{Höhe}\\ \kreuz{Seitenhöhe}\\ \kreuz{Seitenkante}
\teil Pyramide: Grundfläche $27$ cm², Höhe $8$ cm – Volumen? $V$ = \leerfeld[cm³]
\teil Quadratische Pyramide: Grundkante $6$ cm, Höhe $7$ cm – Volumen? $V$ = \leerfeld[cm³]
\end{teile}
\end{aufgabe}

\begin{aufgabe}{Pyramide – Prüfungsaufgaben}
\begin{teile}
\teil Das Dach eines Gartenhauses ist eine quadratische Pyramide mit der Grundkante $4{,}5$ m und der Höhe $1{,}8$ m. Zeichne in das Schrägbild die Höhe des Dachs ein. Beschrifte die Höhe und eine Grundkante mit den Maßen \hfill (P10 2015 OS).

\pyramide{4}{4}{2}{}{}{}
\teil Ein Briefbeschwerer aus Glas ist eine quadratische Pyramide mit der Grundkante $6$ cm und der Höhe $4{,}5$ cm. Zeichne in das Schrägbild die Höhe ein. Beschrifte die Höhe und eine Grundkante mit den Maßen \hfill (P10 2015 OS).

\pyramide{4}{4}{3}{}{}{}
\teil Eine Künstlerin baut eine quadratische Pyramide aus Holzplatten: Grundkante $1{,}20$ m, Höhe $0{,}90$ m. Sie sagt: „Für eine Seitenfläche brauche ich ca. $0{,}65$ m² Holz.“ Weise nach, dass sie recht hat \hfill (P10 2015 OS).
\teil Ein Oberlicht aus Glas ist eine quadratische Pyramide: Grundkante $2{,}40$ m, Höhe $1{,}50$ m. Der Glaser schätzt: „Je Seitenfläche brauche ich ca. $2{,}3$ m² Glas.“ Weise nach, dass die Schätzung stimmt \hfill (P10 2015 OS).
\end{teile}
\end{aufgabe}

% Anlauf (ohne Prüfkennung)
\begin{aufgabe}{Kegel – Anlauf}
\begin{teile}
\teil Die Strecke mit ? am Kegel – Höhe, Mantellinie oder Radius? Kreuze an.\\ \kreuz{Höhe}\\ \kreuz{Mantellinie}\\ \kreuz{Radius}

\kegel{1.5}{3}{}{?}{}
\teil Kegel: Radius $3$ cm, Höhe $7$ cm – Volumen? Runde auf eine Stelle. $V$ ≈ \leerfeld[cm³]
\teil Kegel: Durchmesser $8$ cm, Höhe $9$ cm – Volumen? Runde auf eine Stelle. $V$ ≈ \leerfeld[cm³]
\end{teile}
\end{aufgabe}

\begin{aufgabe}{Kegel – Prüfungsaufgaben}
\begin{teile}
\teil Ein Turm besteht aus einem $18$ m hohen Zylinder mit einem Kegeldach: Radius $3{,}5$ m, Mantellinie $6{,}2$ m. Das Dach bekommt Schiefer für $45$ € je m². Was kostet das, und wie hoch ist das Bauwerk insgesamt \hfill (P10 2026 FOR)? \leerfeld[€] ; \leerfeld[m]
\teil Ein Pavillon hat $6{,}5$ m hohe zylinderförmige Wände und ein Kegeldach: Radius $2{,}4$ m, Mantellinie $3$ m. Das Dach bekommt Kupferblech für $38$ € je m². Was kostet das, und wie hoch ist das Bauwerk insgesamt \hfill (P10 2026 FOR)? \leerfeld[€] ; \leerfeld[m]
\teil Ein Wasserturm besteht aus einem $21$ m hohen Zylinder mit einem Kegeldach: Radius $5{,}5$ m, Mantellinie $7{,}3$ m. Das Dach bekommt Ziegel für $28$ € je m². Was kostet das, und wie hoch ist das Bauwerk insgesamt \hfill (P10 2026 FOR)? \leerfeld[€] ; \leerfeld[m]
\teil Ein Kegel und ein Zylinder sind gleich hoch. Mia behauptet: „Ist der Radius des Kegels sechsmal so groß wie der Radius des Zylinders, hat der Kegel das doppelte Volumen.“ Prüfe die Aussage mit einer Rechnung \hfill (P10 2026 FOR).
\teil Ein Kegel und ein Zylinder sind gleich hoch. Jonas behauptet: „Ist der Radius des Kegels anderthalbmal so groß wie der Radius des Zylinders, hat der Kegel halb so viel Volumen wie der Zylinder.“ Prüfe die Aussage mit einer Rechnung \hfill (P10 2026 FOR).
\end{teile}
\end{aufgabe}

% Anlauf (ohne Prüfkennung)
\begin{aufgabe}{Kugel – Anlauf}
\begin{teile}
\teil Die Strecke mit ? an der Kugel – Radius oder Durchmesser? Kreuze an.\\ \kreuz{Radius}\\ \kreuz{Durchmesser}

\kugel{1.5}{?}
\teil Kugel: Radius $3$ cm – Volumen? Runde auf eine Stelle. $V$ ≈ \leerfeld[cm³]
\teil Kugel: Durchmesser $16$ cm – Volumen? Runde auf eine Stelle. $V$ ≈ \leerfeld[cm³]
\end{teile}
\end{aufgabe}

\begin{aufgabe}{Kugel – Prüfungsaufgaben}
\begin{teile}
\teil Eine Bowlingkugel hat den Durchmesser $21{,}6$ cm. Berechne ihr Volumen mit der Formel aus der Formelsammlung \hfill (P10 2016 OS). $V$ ≈ \leerfeld[cm³]
\teil Ein Wasserball hat den Durchmesser $40$ cm. Berechne sein Volumen mit der Formel aus der Formelsammlung \hfill (P10 2016 OS). $V$ ≈ \leerfeld[cm³]
\teil Eine Boule-Kugel aus Stahl wiegt $700$ g, Stahl hat die Dichte $7{,}85$ g/cm³. Berechne das Volumen der Kugel \hfill (P10 2016 OS). $V$ ≈ \leerfeld[cm³]
\teil Eine große Glasmurmel wiegt $150$ g, Glas hat die Dichte $2{,}5$ g/cm³. Berechne das Volumen der Murmel \hfill (P10 2016 OS). $V$ = \leerfeld[cm³]
\end{teile}
\end{aufgabe}
